\chapter{Configuration and Environment}
\label{chap:appendices}

This appendix provides the exact configuration files, time complexity breakdown, and hardware specifications used to conduct the systematic simulations.

\section{Hardware Specification}
\label{sec:app-hardware}
Due to the computational intensity of generating and evaluating over 1.8 million synthetic datasets, the entire experimental pipeline was executed on a high-performance local server, leveraging parallelized multi-core computing. The exact hardware specifications of the workstation are detailed below:
\begin{itemize}
    \item \textbf{CPU:} AMD Ryzen 9 9900X (12 Cores, 24 Threads)
    \item \textbf{RAM:} 48GB Corsair Vengeance DDR5 6000MHz CL36
    \item \textbf{GPU:} Dual NVIDIA GeForce RTX 3090 (2x 24GB GDDR6X)
    \item \textbf{Storage:} 1TB NVMe SSD (7250 MB/s Sequential Read)
    \item \textbf{OS:} Kubuntu 24.04 LTS
\end{itemize}

\section{Computational Complexity and Parallel Architecture}
\label{sec:app-time-complexity}
To ensure statistical robustness, the regression experiment was run for $M = 1,000$ iterations per scenario across $N \in \{100, 500, 1000\}$ and $p \in \{2, 5, 10\}$. This extensive grid required the synthesis and fitting of 1.8 million individual datasets.

\subsection{Parallel Execution Framework}
Executing such a massive volume of isolated mathematical simulations sequentially is computationally prohibitive. To overcome this, the \texttt{run\_all.sh} master script acts as an advanced orchestrator, distributing the computational load across the 24 available CPU threads. 

The core of this efficiency lies in an optimized \textbf{dynamic work-stealing queue} architecture built for the R and Python subprocesses. Instead of statically assigning datasets to workers (which causes bottlenecking due to the highly variable convergence times of logistic models), a central queue allows idle CPU threads to dynamically pull the next available dataset. Race conditions between concurrent processes are strictly prevented using atomic file renaming operations (\texttt{file.rename}), ensuring that multiple workers do not attempt to overwrite the same regression output.

\subsection{Simulation Runtime}
This highly parallelized architecture achieved near-linear scaling. The time complexity breakdown for the full $M$-run is as follows:
\begin{itemize}
    \item \textbf{Total Worker Time (Sequential Equivalent):} 668,697.38 seconds ($\approx 185.75$ hours)
    \item \textbf{Total Wall-Clock Time (Parallelized):} 37,409.08 seconds ($\approx 10.39$ hours)
    \item \textbf{Linear Regression (Continuous \& Binary):} 1,909.65 wall-clock seconds ($\approx 0.53$ hours)
    \item \textbf{Logistic Regression (Continuous \& Binary):} 35,499.42 wall-clock seconds ($\approx 9.86$ hours)
\end{itemize}

Logistic regression synthesis and fitting represented approximately 95\% of the total computational overhead. This is primarily due to the iterative nature of maximum likelihood estimation and the high rate of non-convergence (complete separation) encountered when attempting to fit complex models on small samples ($N=100$, $p=10$). By leveraging the parallel queue, the actual elapsed time was reduced by a factor of 18, transforming an impossible month-long computation into an overnight run.

\section{Run Configurations (JSON)}
\label{sec:app-json}

The exact boundaries of the simulation space were defined by declarative JSON configurations, guaranteeing reproducibility.

\subsection{Run 1: Clustering Experiment Configuration}
\begin{verbatim}
{
  "simulation": {
    "N": 1000,
    "n": 5,
    "m": 100,
    "random_seed_base": 16
  },
  "parameters": {
    "p": [2, 5, 10],
    "k": [2, 3, 4],
    "separation": [0.1, 2, 6, 10],
    "rho": [0.0, 0.4],
    "distribution": [
      { "name": "normal" },
      { "name": "gamma", "shape": 1.5, "rate": 5.0 }
    ]
  }
}
\end{verbatim}

\subsection{Run 2: Regression Experiment Configuration}
\begin{verbatim}
{
  "simulation": {
    "N": [100, 500, 1000],
    "p": [2, 5, 10],
    "var_type": ["continuous", "binary"],
    "random_seed_base": 16,
    "M": 1000
  },
  "parameters": {
    "rho": [0.0, 0.3, 0.6],
    "sigma_2": [0.5, 1.0, 2.0],
    "p1": [0.5],
    "beta": [1, -1, 1, -1, 1, -1, 1, -1, 1, -1, 1]
  },
  "synthesis": {
    "continuous_methods": ["cart", "norm"],
    "binary_methods": ["cart", "logreg"]
  }
}
\end{verbatim}
